\documentclass[11pt]{article}
\usepackage[margin=1.05in]{geometry}
\usepackage{amsmath,amssymb,amsthm}
\usepackage{url}\urlstyle{tt}
\usepackage{array}

% ---------- theorem environments ----------
\theoremstyle{plain}
\newtheorem{theorem}{Theorem}[section]
\newtheorem{proposition}[theorem]{Proposition}
\newtheorem{lemma}[theorem]{Lemma}
\newtheorem{corollary}[theorem]{Corollary}
\theoremstyle{definition}
\newtheorem{definition}[theorem]{Definition}
\newtheorem{example}[theorem]{Example}
\theoremstyle{remark}
\newtheorem{remark}[theorem]{Remark}

% ---------- notation ----------
\newcommand{\Q}{\mathbb{Q}}
\newcommand{\Z}{\mathbb{Z}}
\newcommand{\R}{\mathbb{R}}
\newcommand{\Mah}{\mathrm{M}}
\newcommand{\Rat}{\mathrm{Rat}}
\DeclareMathOperator{\Arg}{Arg}
\DeclareMathOperator{\Res}{Res}
\DeclareMathOperator{\lcm}{lcm}
\newcommand{\tors}{\Q/\Z}

% ---------- epistemic tags (project convention) ----------
\newcommand{\F}{\textnormal{[forced]}}
\newcommand{\C}{\textnormal{[computed]}}
\newcommand{\Pl}{\textnormal{[plausible]}}
\newcommand{\Op}{\textnormal{[open]}}

\title{Relational Charge on the Spectral Semiring:\\
Gauge Rigidity, Coherence Types, and a Reference-Free Parity Floor}
\author{math-research-pipelines technical note}
\date{1 July 2026}

\begin{document}
\maketitle

\begin{abstract}
Which part of an algebraic integer's phase structure belongs to the number
itself, and which only to a chosen reference ray? We refactor the emission
algebra's charge into a purely relational invariant --- the pair charge
$t_{\mathrm{rel}}(\alpha,\beta)=\frac{1}{2\pi}(\Arg\alpha-\Arg\beta)
\in\Q/\Z$, defined when rational --- and determine exactly how much
survives. The relational data is a groupoid cocycle whose trivializations
are the absolute charges; the reference ray is a gauge choice. Rigidity:
conjugation-closed spectra force every internally rational object to be
absolutely admissible, with anchor $n\in\{m,2m\}$ over the relational group
$\Z/m\Z$ --- the absolute theory is the canonical section --- and the
composition laws descend verbatim. On the inadmissible sector the theory is
genuinely finer: objects carry a computable coherence type, and a
\emph{modulus-pinning theorem} --- an irreducible polynomial with a
uniquely attained root modulus admits no root-of-unity ratio between
distinct roots --- proves every Salem number relationally inert and no two
Salem numbers circle-locked, while irreducible non-inert inadmissible
objects do exist (twisted shells $q(x^k)$).

All of this is executed exactly: a shell-complete, float-free contact
procedure corroborates the pinning theorem on forty-six certified
instances, including an exhaustive degree-$12$ small-coefficient census
($37/37$ inert), cross-checked in a second computer-algebra stack. Last ---
and conditionally on the companion note's emission gap, the one hypothesis
not proved here beyond degree two --- the parity-graded Mahler floor is
evaluated reference-free: the parity bit is the golden pair's internal
order-$2$ relation $t_{\mathrm{rel}}(\varphi,\varphi')=\tfrac12$, a sign
twist absorbs the second anchor sector, and
$\mu_{\mathrm{rel}}=\mu$; that identity and the attainments are
unconditional, only the numerical floor values rest on the gap.
Unconditional proofs, gap-conditional results, and per-instance
computations are tagged apart throughout.
\end{abstract}

\section{Introduction}

\subsection{The question}

The emission algebra of the companion note~\cite{CMC} grades its objects by two
orthogonal readouts: a magnitude (the Mahler measure $\Mah$, subject to the
emission gap $\Mah\in\{1\}\cup[\varphi,\infty)$ with $\varphi=\frac{1+\sqrt5}{2}$)
and a phase (a charge $\chi$ valued in a finite cyclic group $\Z/n\Z$, defined when
every conjugate's argument lies on a finite rational-angle lattice). Both are
\emph{object-attached}: the charge of a conjugate $\alpha$ is read against the
fixed reference ray $\Arg=0$, and an object either possesses a charge group or ---
as for every Salem number~\cite{Salem,Bertin} --- possesses none.

This note asks what changes when the phase character is attached instead to the
\emph{relationship between oscillators}. Each conjugate $\alpha=|\alpha|\,
e^{2\pi i\,t(\alpha)}$ is an oscillator with growth $|\alpha|$ and phase
$t(\alpha)\in\R/\Z$; the relational datum is the phase \emph{difference}
\[
t_{\mathrm{rel}}(\alpha,\beta)\;=\;t(\alpha)-t(\beta)\;\in\;\tors ,
\]
declared whenever the difference is rational, with no reference ray anywhere. Three
questions organize the study:
\begin{enumerate}
\item[(i)] \textbf{Gauge.} In what precise sense is the absolute charge a
  gauge-fixing of the relational one, and what plays the role of the gauge group?
\item[(ii)] \textbf{Rigidity.} Does the relational theory admit objects the
  absolute theory rejects --- oscillator systems internally in tune but globally
  out of tune with the reference ray?
\item[(iii)] \textbf{Content.} Which theorems of the absolute theory --- the
  $\lcm$ composition law, the Adams/CRT projectors, the parity-graded floor
  $\mu(n)\in\{\varphi,2\}$ --- are secretly relational, and what genuinely new
  invariants appear?
\end{enumerate}

\subsection{Contributions}

\begin{itemize}
\item \textbf{The gauge structure} (\S\ref{sec:relational}). The relational charge
  is a $\tors$-valued cocycle on a coherence groupoid; absolute charge is a
  trivialization; the reference ray is the gauge; a global rotation
  $\alpha\mapsto e^{2\pi ic}\alpha$ shifts every absolute phase by $c$ and fixes
  every relative one. What the language buys is the exact content of
  \S\ref{sec:rigidity}: rigidity says the algebraic category admits
  essentially one section --- integrality fixes the gauge. \F
\item \textbf{Rigidity} (\S\ref{sec:rigidity}). Integrality fixes the gauge: for a
  conjugation-closed spectrum, all pairwise differences rational forces all phases
  rational. Relational admissibility therefore coincides with absolute
  admissibility, with the absolute group pinned to an anchor $n\in\{m,2m\}$ over
  the relational group $\Z/m\Z$; both anchors occur, and the relational group can
  be a proper subgroup of the absolute one (witness: $x^4+5x^2+5$, absolute
  $\Z/4\Z$, relational $\Z/2\Z$, $\Mah=5$). Read this as a canonicity
  theorem, not a redundancy: it proves the absolute framework of~\cite{CMC}
  is the unique natural presentation of the relational one on its own domain,
  while the relational language retains purchase exactly where absolute
  methods have none --- the inadmissible sector. \F
\item \textbf{Descent of the composition laws} (\S\ref{sec:descent}). Coherence
  classes are preserved by the Adams operations, which multiply relative charges by
  $k$; the tensor law $\Delta(A\otimes B)\supseteq\Delta(A)+\Delta(B)$ holds with
  equality on admissible inputs (the $\lcm$ law, relationally), while strict
  growth occurs exactly through \emph{offset cancellation} --- the tensor product
  internalizes cross-object phase relations. \F
\item \textbf{A reference-free parity floor} (\S\ref{sec:parity}). The parity bit
  of the charge--measure coupling is intrinsic: it is the realized difference
  $t_{\mathrm{rel}}(\varphi,\varphi')=\tfrac12$ (since $\varphi/\varphi'=
  -\varphi^2$), not the lattice membership of the ray $\pi$. A sign-twist
  involution $p(x)\mapsto(-1)^{\deg p}p(-x)$ exchanges the two anchor sectors at
  odd $m$ and preserves the Mahler measure, so the relational floor equals the
  absolute one: $\mu_{\mathrm{rel}}(m)=\varphi$ for even $m$ (attained), and
  $\mu_{\mathrm{rel}}(m)=\mu(m)$ for odd $m$. Identity and attainment \F;
  the floor \emph{values} are conditional per Remark~\ref{rem:gapstatus}
\item \textbf{Coherence types of inadmissible objects} (\S\ref{sec:inert}). On the
  sector the absolute theory discards, the relational theory is genuinely finer:
  every object carries a finite coherence type, computable exactly on
  equal-modulus shells by a ratio-resultant / cyclotomic-contact procedure
  with a provable completeness bound; and a \emph{modulus-pinning theorem}
  --- irreducibility plus a uniquely attained root modulus forbids torsion
  ratios between distinct roots --- proves \emph{every} Salem number
  relationally inert \F, corroborated on forty-six certified instances
  (degrees $4$--$10$ including $\beta_4$ and Lehmer's polynomial, an
  exhaustive degree-$12$ small-coefficient census of $37$, and five
  larger-coefficient spot checks). Irreducible non-inert inadmissible
  objects exist: twisted shells $q(x^k)$, witness $x^4+x^2+2$.
\item \textbf{Cross-object locking} (\S\ref{sec:locking}). Two objects can be
  phase-locked while each is internally inert; the notion is decidable on circle
  shells by a mixed ratio resultant. The pinning theorem's mixed form proves
  no two distinct Salem numbers circle-lock \F; the six pairwise scans among
  the canonical instances stand as its exact corroboration.
\end{itemize}

\subsection{Relation to the companion note}

This note is a refactoring and an extension of~\cite{CMC}, not a replacement. Every
statement of~\cite{CMC} about admissible objects is recovered here, usually in a
sharper, reference-free form; the new content lives (a) in the gauge/rigidity
analysis, which explains \emph{why} the absolute theory loses nothing on its own
domain, and (b) on the inadmissible sector, where the absolute theory is silent by
construction and the relational theory supports nontrivial, decidable invariants.
Notation follows~\cite{CMC} throughout. Because~\cite{CMC} is a technical
note of the same project rather than an archived publication, this note is
organized so that its verification burden is explicit. The unconditional
results (\S\S\ref{sec:relational}--\ref{sec:descent} and
\S\S\ref{sec:inert}--\ref{sec:locking}) are fully auditable within this
document. The conditional results (the floor \emph{values} in
Theorem~\ref{thm:parityrel}) are auditable here in the sense that their
dependence on the named hypothesis \textup{(EG)} is explicit and traceable
(Remark~\ref{rem:gapstatus}, Appendix~\ref{app:inherited}); verifying
\textup{(EG)} itself beyond the quadratic case reproduced in
Appendix~\ref{app:quadEG} requires~\cite{CMC}, archived at DOI \url{10.5281/zenodo.21121863}
(within the project repository deposit, v1.0.0).

\medskip
\noindent\textbf{Relation to the wider literature.} A rational phase
difference between conjugates is, equivalently, a torsion multiplicative
relation $\alpha_i=\zeta\,\alpha_j$ between conjugates. Multiplicative
relations connecting conjugate algebraic numbers are a classical subject
(Smyth~\cite{Smyth1986} and the literature following it); the present note
can be read as isolating their torsion part, packaging it as a
groupoid-valued invariant, and supplying a shell-complete exact decision
procedure for it. On the Salem question specifically, the shell-torsion
half of inertness is exactly \emph{non-degeneracy} in the sense of linear
recurrence sequences~\cite{EPSW}; Remark~\ref{rem:degen} records the
bridge, the modulus-pinning proof, and what the coherence type adds beyond
it. The census infrastructure connects to the small-Salem
tables originating with Boyd~\cite{Boyd1977} and extended by
Mossinghoff~\cite{Mossinghoff}; the unconditional floor landscape below the
gap is Dobrowolski's bound~\cite{Dobrowolski}. We are not aware of a prior
treatment of the coherence type as an invariant of charge-inadmissible
objects; the project's standing citation-provenance discipline applies (see
the provenance block following the references).

\subsection{A note on rigor}

Every numerical assertion below was verified in exact arithmetic
(\S\ref{sec:method}); in this note the entire verification path is integer and
rational polynomial arithmetic --- no floating-point number appears anywhere, in
decisions or in display. Each result carries an epistemic tag: \F{} (proved),
\C{} (verified exhaustively over a stated finite range), \Pl{} (strongly
evidenced, not proved in full generality)\footnote{That is, \Pl{} marks a
conjecture with a specified evidence level: supported by exhaustive
computation over stated ranges and by structural argument, but not proved.
The tag system is project-internal; this footnote fixes its dictionary.}, \Op{} (open). Statements inheriting a
caveated status from~\cite{CMC} say so explicitly, and
Appendix~\ref{app:inherited} lists every statement inherited
from~\cite{CMC} with the status tagged there and the place it enters this
note. The conditional content of this note is confined to the two floor
evaluations of Theorem~\ref{thm:parityrel} (Remark~\ref{rem:robust});
every structural statement --- rigidity, gauge, descent, coherence types,
locking --- is unconditional. We do not claim to resolve
Lehmer's problem or the general no-Salem statement; we characterize their
interaction with the relational phase structure.

\section{Preliminaries}\label{sec:prelim}

\subsection{Objects and operators}

Fix the ring $\overline{\Z}$ of algebraic integers. An \emph{object} is a finite
multiset $O=\{\alpha_1,\dots,\alpha_n\}$ of nonzero algebraic integers that is
closed under complex conjugation --- equivalently, the root multiset of a monic
$p\in\Z[x]$ with $p(0)\neq0$ (not necessarily irreducible), equivalently the
spectrum of an integer matrix. The algebra's operations are those of~\cite{CMC}:
direct sum $\oplus$ (multiset union; on polynomials, the product), tensor
$\otimes$ (eigenvalue products; Kronecker product of matrices), and the Adams
operations $\psi^k$ ($\lambda\mapsto\lambda^k$ on each eigenvalue), $k\ge1$.
The semiring/$\lambda$-ring structure of these operations --- associativity,
distributivity of $\otimes$ over $\oplus$, identities, and
$\psi^a\circ\psi^b=\psi^{ab}$ --- is established in~\cite{CMC} and is not
re-verified here; this note uses only the operations and their effect on
angles, $t(\alpha\beta)=t(\alpha)+t(\beta)$ and $t(\lambda^k)=k\,t(\lambda)$,
not the axiomatic package.

\subsection{The angle coordinate and the absolute charge}

\begin{definition}[Angle]\label{def:angle}
For $\alpha\in\mathbb{C}^\times$ write $t(\alpha)=\frac{1}{2\pi}\Arg\alpha
\in\R/\Z$, so $\alpha=|\alpha|\,e^{2\pi i\,t(\alpha)}$.
\end{definition}

Because $p$ has real coefficients, the angle multiset $t(O)$ is closed under
negation in $\R/\Z$; we use this constantly.

\begin{definition}[Absolute charge, {\cite[Def.~2.4]{CMC}}]\label{def:absolute}
An object is \emph{charge-admissible} if every conjugate $\alpha$ satisfies
$\alpha^n\in\R_{>0}$ for some $n\ge1$, i.e.\ $t(\alpha)\in\tors$ for every
conjugate. The \emph{charge group} is $\Z/n\Z$ for $n$ the least common
denominator of the $t(\alpha)$, and the charge of $\alpha$ is
$\chi_n(\alpha)=n\,t(\alpha)\bmod n$.
\end{definition}

The magnitude character is the Mahler measure~\cite{Mahler}
$\Mah(O)=\prod_i\max(1,|\alpha_i|)$, with Kronecker's theorem ($\Mah=1$ iff every
root is a root of unity) and the emission gap
$\Mah\in\{1\}\cup[\varphi,\infty)$ carrying the statuses recorded
in~\cite[\S2, \S8]{CMC}. We write $\mu(n)$ for the absolute floor of~\cite[Thm.~5.1]{CMC}:
the infimum of $\Mah>1$ over charge-admissible objects with charge group $\Z/n\Z$.

\begin{remark}[Status of the emission gap, and the classical floors]
\label{rem:gapstatus}
Three floor statements are in play and must not be conflated.
(i)~Kronecker's characterization of $\Mah=1$ is classical and
proven~\cite{Kronecker}. (ii)~The unconditional classical lower bound is
Smyth's theorem~\cite{Smyth}: a \emph{non-reciprocal} monic integer
polynomial has $\Mah\ge\theta_0=1.32471\ldots$, the real root of $x^3-x-1$
(the smallest Pisot number, often called the plastic number; we refer to
$\theta_0$ as Smyth's constant); for reciprocal polynomials the existence of any uniform bound $>1$ is
Lehmer's problem~\cite{Lehmer}, open. (iii)~The emission gap used here ---
$\Mah\in\{1\}\cup[\varphi,\infty)$ \emph{for charge-admissible objects} ---
is \cite[Prop.~2.3]{CMC}, treated here as a named hypothesis \textup{(EG)}
(its quadratic case is elementary and reproduced self-containedly in
Appendix~\ref{app:quadEG}):
\F{} for the integer-quadratic spectrum, \C/\Pl{} in general. Since $\varphi>\theta_0$, statement~(iii) is not implied by
Smyth's theorem even on the non-reciprocal part; nor is it Lehmer's
conjecture, being restricted to the rational-angle sector. In interval
terms: Smyth's theorem already excludes $(1,\theta_0)$ for every
non-reciprocal polynomial; \textup{(EG)} additionally excludes
$[\theta_0,\varphi)$ on the charge-admissible sector, where it binds
reciprocal and non-reciprocal objects alike. Consequently,
wherever this note evaluates a floor \emph{as} $\varphi$, the attainment is
\F{} while the identification of the infimum as $\varphi$ (rather than some
value in $[\theta_0,\varphi)$) is conditional on \textup{(EG)};
Theorem~\ref{thm:parityrel} and the table of \S\ref{sec:open} carry this
tag explicitly.
\end{remark}

\section{The relational charge}\label{sec:relational}

\subsection{Coherence, classes, and the relational group}

\begin{definition}[Pair charge and coherence]\label{def:pair}
Conjugates $\alpha,\beta\in O$ are \emph{coherent}, written $\alpha\approx\beta$,
if $t(\alpha)-t(\beta)\in\tors$; in that case their \emph{relative charge} is
\[
t_{\mathrm{rel}}(\alpha,\beta)\;=\;t(\alpha)-t(\beta)\;\in\;\tors .
\]
\end{definition}

Coherence is an equivalence relation ($t(\alpha)-t(\alpha)=0$; symmetry by
negation; transitivity because rationals are closed under addition), so $O$
partitions into \emph{coherence classes} $C_1,\dots,C_r$.

\begin{definition}[Class groups, relational group, coherence type]\label{def:type}
For a class $C$, its \emph{class group} is the subgroup
$\Delta(C)=\langle\,t_{\mathrm{rel}}(\alpha,\beta):\alpha,\beta\in
C\,\rangle\le\tors$. The \emph{relational charge group} of $O$ is
$\Delta(O)=\sum_i\Delta(C_i)\le\tors$. The object is \emph{relationally
admissible} if it has a single coherence class. The \emph{coherence type} of $O$
is the triple: the partition $\{C_i\}$, the class groups $\{\Delta(C_i)\}$, and
the involution that negation induces on the classes
(Theorem~\ref{thm:structure}).
\end{definition}

\begin{example}\label{ex:firsttype}
For $x^3-2$ the angles are $\{0,\tfrac13,\tfrac23\}$: every pairwise
difference is a multiple of $\tfrac13$, so there is a single class with
$\Delta=\Z/3\Z$, fixed by the negation involution. For $x^4+5x^2+5$
(Example~\ref{ex:drop}) the angles are $\{\tfrac14,\tfrac34\}$: again a
single class, now with $\Delta=\Z/2\Z$. For a Salem quartic such as
$\beta_4$ the type is computed in \S\ref{sec:inert}: rational block
$\{\tau,1/\tau\}$ plus two mirrored circle singletons.
\end{example}

\begin{proposition}[Cyclicity, relational form]\label{prop:cyclic}
Every class group and every relational group is a finite cyclic group $\Z/m\Z$;
no non-cyclic group occurs. \F
\end{proposition}

\begin{proof}
Each is a subgroup of $\tors$ generated by finitely many rationals, hence a finite
subgroup of $\tors$; $\tors$ has exactly one subgroup of each finite order,
namely $\langle 1/m\rangle\cong\Z/m\Z$. This is~\cite[Thm.~3.1]{CMC} with the
generators taken to be differences rather than absolute angles.
\end{proof}

\subsection{The cocycle, the trivialization, and the gauge}

\begin{proposition}[Cocycle and trivialization]\label{prop:gauge}
\leavevmode
\begin{enumerate}
\item[(i)] On coherent triples, $t_{\mathrm{rel}}(\alpha,\beta)+
  t_{\mathrm{rel}}(\beta,\gamma)=t_{\mathrm{rel}}(\alpha,\gamma)$: the relative
  charge is a $\tors$-valued cocycle on the coherence groupoid (objects: the
  conjugates; morphisms: coherent pairs).
\item[(ii)] If $O$ is charge-admissible with charge group $\Z/n\Z$, then on every
  pair $t_{\mathrm{rel}}(\alpha,\beta)=\frac{1}{n}\bigl(\chi_n(\alpha)-
  \chi_n(\beta)\bigr)$: the relational charge is the coboundary of the absolute
  one.
\item[(iii)] Any two functions $s:C\to\tors$ with $s(\alpha)-s(\beta)=
  t_{\mathrm{rel}}(\alpha,\beta)$ on a class $C$ differ by a constant.
\item[(iv)] The rotation $\rho_c:\alpha\mapsto e^{2\pi ic}\alpha$ (an operation on
  the ambient analytic picture; for irrational $c$ it leaves the category of
  algebraic integers) shifts every absolute angle by $c$ and fixes every relative
  charge. \F
\end{enumerate}
\end{proposition}

\begin{proof}
(i) and (iii) are immediate from Definition~\ref{def:pair}; (ii) is
Definition~\ref{def:absolute}; (iv) is $t(\rho_c\alpha)=t(\alpha)+c$.
\end{proof}

We call the choice of the reference ray $\Arg=0$ --- equivalently, the choice of
constant in (iii) --- the \emph{gauge}. The absolute theory of~\cite{CMC} is the
relational theory plus this gauge choice. The terminology is structural, not
physical: the circle acts as a structure group on phase assignments,
Proposition~\ref{prop:gauge}(iii) exhibits the trivializations as a torsor
under it, and integrality (\S\ref{sec:rigidity}) provides the canonical
section; no connection to physical gauge theory is implied. The cocycle
language earns its keep twice over: it is the functorial carrier for the
descent analysis of \S\ref{sec:descent}, and it reduces the proof of
Theorem~\ref{thm:rigidity} to a single coset computation. The next section
shows that on the
algebra's own objects the choice is almost costless: integrality fixes the gauge.

\section{Rigidity: integrality fixes the gauge}\label{sec:rigidity}

\begin{theorem}[Rigidity]\label{thm:rigidity}
Let $O$ be an object (conjugation-closed) that is relationally admissible with
relational group $\Delta(O)=\Z/m\Z$. Then every angle $t(\alpha)$ is rational:
$O$ is charge-admissible, and its absolute charge group $\Z/n\Z$ satisfies
\[
n\in\{m,\;2m\}.
\]
Both anchors occur. \F
\end{theorem}

\begin{proof}
All pairwise differences lie in $\frac1m\Z/\Z$. Fix any conjugate $\alpha$.
Since $O$ is closed under complex conjugation, $\beta_0:=\overline{\alpha}$
is among the conjugates, with $t(\beta_0)=-t(\alpha)$; since $O$ is relationally admissible --- a single coherence class ---
every pair of conjugates is coherent, in particular $\alpha\approx\beta_0$, so
$t(\alpha)-(-t(\alpha))=2\,t(\alpha)\in\frac1m\Z/\Z$, whence
$t(\alpha)\in\frac{1}{2m}\Z/\Z$. Thus all angles are rational and lie in
$\frac{1}{2m}\Z/\Z$, so the least common denominator $n$ divides $2m$. And $m$ divides $n$:
every angle lies in $\frac1n\Z/\Z$, hence every difference does; the
differences generate $\frac1m\Z/\Z$, so
$\frac1m\Z/\Z\subseteq\frac1n\Z/\Z$, i.e.\ $m\mid n$. Hence
$n\in\{m,2m\}$. For the anchors:
$x^m-2$ has angles $\{j/m\}$, so $n=m=\Delta$; $x^3+2$ has angles
$\{\tfrac16,\tfrac12,\tfrac56\}$, so $n=6$ while the differences generate
$\tfrac13\Z/\Z$, i.e.\ $m=3$, $n=2m$ (ledger entry~B).
\end{proof}

\begin{corollary}\label{cor:nonew}
Relational admissibility coincides with absolute admissibility on the objects of
the algebra. In particular the relational theory admits no new objects: an
oscillator system that is internally in tune (all pairwise phase differences
rational) is automatically in tune with the reference ray. \F
\end{corollary}

The mechanism deserves emphasis: it is the \emph{integrality} of the object ---
real coefficients, hence conjugation-closed spectra --- that pins the gauge. The
rotations $\rho_c$ of Proposition~\ref{prop:gauge}(iv) act freely on the analytic
envelope, but the algebraic locus meets each orbit only at rational offsets.
Rigidity is thus a feature, not a limitation: it is the uniqueness half of a
gauge-fixing statement --- the absolute theory of~\cite{CMC} is the canonical
section --- and it locates the relational theory's genuine domain, the
inadmissible sector of \S\ref{sec:inert}, where there is no section to fix.
Equivalently, rigidity is a \emph{completeness guarantee}: no information is
lost in the relational reformulation --- the prerequisite for trusting it on
the sector where it extends reach.

\begin{theorem}[Structure of partial coherence]\label{thm:structure}
Negation $t\mapsto -t$ permutes the coherence classes of any object. A class $C$
with $-C=C$ has all angles rational ($t(C)\subset\tors$); a class with
$-C\neq C$ carries an offset: $t(C)\subseteq c+\frac{1}{m_C}\Z/\Z$ for a class
constant $c$ that is irrational (well defined modulo $\frac{1}{m_C}\Z/\Z$), and
$t(-C)=-t(C)$ carries the offset $-c$. Every object therefore decomposes into a
\emph{rational block} (the union of the self-paired classes, which is
charge-admissible on its own) and mirrored pairs of irrationally offset classes.
\F
\end{theorem}

\begin{proof}
Negation preserves rationality of differences, hence maps classes to classes. If
$-C=C$ then for $\alpha\in C$ the angle $-t(\alpha)$ belongs to $t(C)$, so
$2t(\alpha)\in\Delta(C)\subset\tors$ and $t(\alpha)$ is rational; conversely all
differences inside a class being rational, one rational member makes all members
rational. If $-C\neq C$, pick $\alpha_0\in C$ and set $c=t(\alpha_0)$; every
member differs from $\alpha_0$ by an element of $\frac{1}{m_C}\Z/\Z$, giving the
coset form, and $c$ must be irrational or else $C$ would be rational and
self-paired by the previous sentence.
\end{proof}

\begin{lemma}[Odd anchors are full]\label{lem:oddfull}
If $O$ is charge-admissible with odd absolute group $\Z/n\Z$, then
$\Delta(O)=\Z/n\Z$: at odd order the relational and absolute groups coincide. \F
\end{lemma}

\begin{proof}
Say $\Delta(O)=\Z/d\Z$ with $d\mid n$. All angles lie in a single coset
$t_0+\frac1d\Z/\Z$ (differences generate $\frac1d\Z/\Z$ and the object is one
class since all angles are rational). Negation closure forces
$2t_0\in\frac1d\Z/\Z$, so $t_0\in\frac{1}{2d}\Z/\Z$ and the angles lie in
$\frac{1}{2d}\Z/\Z$; hence $n\mid 2d$. With $n$ odd this gives $n\mid d$, so
$d=n$.
\end{proof}

\begin{example}[The group drop, and its absence at odd order]\label{ex:drop}
$g=x^4+5x^2+5$ is irreducible (Eisenstein at $5$); its $x^2$-roots
$\frac{-5\pm\sqrt5}{2}$ are both negative, so all four roots are purely
imaginary: angles $\{\tfrac14,\tfrac34\}$, absolute group $\Z/4\Z$, differences
$\{0,\tfrac12\}$, relational group $\Z/2\Z$, and $\Mah(g)=5$ (the squared
moduli are $\frac{5\mp\sqrt5}{2}\approx1.38$ and $3.62$, both exceeding $1$
exactly because $5-\sqrt5>2\iff\sqrt5<3$; hence $\Mah=|g(0)|=5$). The absolute $n$ thus conflates two invariants ---
the relational group and the anchor --- which Lemma~\ref{lem:oddfull} shows can
differ only at even $n$. (Ledger entry~E; compare the catalog seed
$K=x^4+5x^2-5$ of~\cite{CMC}, whose sign-mirror this is: $K$ keeps the full
$\Delta=\Z/4\Z$ because its real pair sits at angles $\{0,\tfrac12\}$.) \F
\end{example}

\begin{corollary}[Salem geometry, relational form]\label{cor:salemgeom}
Let $\tau$ be a Salem number of degree $2s+2$, with conjugates $\tau$, $1/\tau$,
and $s$ unit-circle pairs $\lambda_j,\overline{\lambda_j}$. Then:
\begin{enumerate}
\item[(i)] every circle conjugate has irrational angle (a unimodular number with
  rational angle is a root of unity, contradicting~\cite[Lem.~6.1]{CMC});
\item[(ii)] no circle conjugate is coherent with its own complex conjugate
  (their difference is $2t(\lambda_j)$, irrational), so circle classes come in
  mirrored pairs;
\item[(iii)] the rational block is exactly $\{\tau,1/\tau\}$, with trivial class
  group (both angles are $0$);
\item[(iv)] the only undetermined coherence is \emph{within} a half-block
  $\{\lambda_1,\dots,\lambda_s\}$ of circle conjugates --- whether some
  $t(\lambda_i)-t(\lambda_j)$ is rational while both angles are irrational.
\end{enumerate}
\F{} for (i)--(iii); (iv) is the question \S\ref{sec:inert} decides per instance.
\end{corollary}

\section{Descent of the composition laws}\label{sec:descent}

\begin{proposition}[Adams operations]\label{prop:adams}
For $k\ge1$, the map $\psi^k$ preserves the coherence partition (as a bijection of
index sets) and multiplies every relative charge by $k$:
$t_{\mathrm{rel}}(\alpha^k,\beta^k)=k\,t_{\mathrm{rel}}(\alpha,\beta)$, so
$\Delta(\psi^kO)=k\,\Delta(O)$. In particular the CRT/Adams projectors
of~\cite[Thm.~3.3]{CMC} descend: for $\Delta(O)=\Z/m\Z$ with
$m=\prod p_i^{e_i}$, the operator $\psi^{m/p_i^{e_i}}$ projects the relational
group onto its $p_i$-primary component. \F
\end{proposition}

\begin{proof}
$t(\alpha^k)-t(\beta^k)=k\bigl(t(\alpha)-t(\beta)\bigr)$ in $\R/\Z$, and for an
integer $k\ge1$ this is rational if and only if $t(\alpha)-t(\beta)$ is rational;
so coherence is preserved and reflected, and the class map is the identity on
indices. The group identity and the CRT statement follow as
in~\cite[Thm.~3.3]{CMC}, applied to differences.
\end{proof}

\begin{proposition}[Tensor: containment, the $\lcm$ law, and internalization]
\label{prop:tensor}
For objects $A,B$:
\begin{enumerate}
\item[(i)] a pair $\bigl(\alpha\beta,\;\alpha'\beta'\bigr)$ in $A\otimes B$ is
coherent iff
$\bigl(t(\alpha)-t(\alpha')\bigr)+\bigl(t(\beta)-t(\beta')\bigr)\in\tors$;
consequently $\Delta(A\otimes B)\supseteq\Delta(A)+\Delta(B)$.
\item[(ii)] If $A$ and $B$ are relationally admissible with groups
$\Z/m\Z,\Z/k\Z$, then $A\otimes B$ is relationally admissible with
$\Delta(A\otimes B)=\frac1m\Z+\frac1k\Z=\frac{1}{\lcm(m,k)}\Z$ modulo $\Z$: the
$\lcm$ law of~\cite[Thm.~4.2]{CMC} is a relational statement, and safe
composition $\Leftrightarrow$ coprime orders descends with it
(\emph{safe}: the two factor charges remain independently recoverable from
the composite, \cite[Thm.~4.3]{CMC}).
\item[(iii)] Strict containment in (i) occurs exactly through \emph{offset
cancellation}: classes $C_A\subseteq A$, $C_B\subseteq B$ with irrational offsets
$c_A,c_B$ (Theorem~\ref{thm:structure}) whose sum $c_A+c_B$ is rational make all
of $C_A\cdot C_B$ mutually coherent. Canonically, for any object $O$ the square
$O\otimes O$ contains, for each mirrored pair $(C,-C)$, the rational block
$C\cdot(-C)$ --- in particular the diagonal products
$\alpha\overline{\alpha}=|\alpha|^2$ at angle $0$. The tensor product
internalizes cross-object phase relations, even between relationally inert
factors. \F
\end{enumerate}
\end{proposition}

\begin{proof}
(i) is $t(\alpha\beta)=t(\alpha)+t(\beta)$; the containment takes
$\beta=\beta'$ or $\alpha=\alpha'$. (ii): by rigidity
(Theorem~\ref{thm:rigidity}) admissible inputs have all angles rational, so every
sum in (i) is rational and $A\otimes B$ is one class. For the group: fixing
$\beta=\beta'$ realizes every element of $\Delta(A)=\frac1m\Z/\Z$ as a
difference in $A\otimes B$, and fixing $\alpha=\alpha'$ realizes
$\Delta(B)=\frac1k\Z/\Z$; conversely, by (i) every difference is a sum of
one element of each. Hence $\Delta(A\otimes B)=\frac1m\Z+\frac1k\Z$ modulo
$\Z$, which is $\frac{d}{mk}\Z/\Z=\frac{1}{\lcm(m,k)}\Z/\Z$ with
$d=\gcd(m,k)$. Recoverability at coprime orders is the CRT projection,
realized inside the algebra by the Adams operators
(Proposition~\ref{prop:adams}); at $\gcd(m,k)=d>1$ the surjection
$\Z/m\Z\times\Z/k\Z\twoheadrightarrow\Z/\lcm(m,k)\Z$, given by
$(a,b)\mapsto a\tfrac{L}{m}+b\tfrac{L}{k}\bmod L$ with $L=\lcm(m,k)$, has
kernel cyclic of order $d$, generated by $(m/d,\,-k/d)$
(cf.~\cite[Thm.~4.3]{CMC}).
(iii): if $c_A+c_B\in\Q$ then any two products from $C_A\times C_B$ differ by an
element of $\frac{1}{m_{C_A}}\Z+\frac{1}{m_{C_B}}\Z\subset\tors$; the mirrored
pair of one object supplies $c+(-c)=0$.
\end{proof}

\begin{proposition}[Direct sum and phase locking]\label{prop:locking}
The coherence classes of $A\oplus B$ are generated by those of $A$ and $B$
together with the \emph{cross edges} $\alpha\approx\beta$
($\alpha\in A,\beta\in B$). If one cross pair between classes $C_A, C_B$ is
coherent then all of $C_A\times C_B$ is, and this happens iff the offsets satisfy
$c_A-c_B\in\Q$ (rational blocks having offset $0$). We call $A$ and $B$
\emph{locked along $(C_A,C_B)$} in that case, and \emph{circle-locked} when both
classes consist of unimodular conjugates. Locking between rational blocks is
automatic and carries no information; locking between irrationally offset classes
is a genuine cross-object relational charge, invisible to the absolute theory on
either factor. \F
\end{proposition}

\begin{proof}
Transitivity through the class relations gives the all-or-nothing statement:
$t(\alpha')-t(\beta')=\bigl(t(\alpha')-t(\alpha)\bigr)+
\bigl(t(\alpha)-t(\beta)\bigr)+\bigl(t(\beta)-t(\beta')\bigr)$, a sum of
rationals. The offset criterion restates the definition of the class constants.
\end{proof}

Propositions~\ref{prop:tensor}(iii) and~\ref{prop:locking} are two views of one
phenomenon: $\oplus$ records a lock as a merged class; $\otimes$ records it as a
created rational block.

\section{The parity floor, reference-free}\label{sec:parity}

\begin{lemma}[The golden internal relation]\label{lem:golden}
The golden pair carries an intrinsic order-$2$ relational charge:
\[
\frac{\varphi}{\varphi'}=-\varphi^2,\qquad\text{hence}\qquad
t_{\mathrm{rel}}(\varphi,\varphi')=\tfrac12 .
\]
\F{} (ledger entry~J)
\end{lemma}

\begin{proof}
$\varphi'=-1/\varphi$, so $\varphi/\varphi'=-\varphi^2$, a negative real:
relative angle $\pi$, relative charge $\tfrac12$.
\end{proof}

This is~\cite[Lem.~5.2]{CMC} with the reference ray deleted: what the
$\pi$-ray obstruction tracks is not the lattice membership of an absolute ray but
the realized \emph{difference} $\tfrac12$ inside the golden pair. The
floor-attaining construction of~\cite[Thm.~5.1]{CMC} always realizes it:

\begin{lemma}[$q_k$ realizes the difference $\tfrac12$]\label{lem:qk}
For every $k\ge1$, $q_k=x^{2k}+x^k-1$ has a real root of each sign
(for odd $k$: $(1/\varphi)^{1/k}>0$ and $-\varphi^{1/k}<0$; for even $k$:
$\pm(1/\varphi)^{1/k}$), hence realizes
$t_{\mathrm{rel}}=\tfrac12$; its angle set is all of $\frac{1}{2k}\Z/\Z$, so
$q_k$ is relationally admissible with $\Delta(q_k)=\Z/2k\Z$ (and anchor $n=2k$).
\F
\end{lemma}

\begin{proof}
The $x^k$-substitution gives $y^2+y-1$ with roots $1/\varphi>0$ and $-\varphi<0$
(ledger entry~K); the stated real roots follow by parity of $k$, and the full
angle set is computed in~\cite[Thm.~5.1]{CMC}: inert roots at angles $j/k$,
measure-bearing roots at $(2j+1)/2k$. Differences include $\frac{1}{2k}$, so the
class group is full.
\end{proof}

\begin{lemma}[The sign twist]\label{lem:twist}
The involution $T(p)(x)=(-1)^{\deg p}\,p(-x)$ on objects multiplies every root by
$-1$; it therefore shifts every angle by $\tfrac12$, fixes every difference
(hence the coherence type and $\Delta$), preserves the Mahler measure, and
preserves monic integrality. For odd $m$ it exchanges the two anchor sectors of
Theorem~\ref{thm:rigidity}: objects with $(\Delta,n)=(\Z/m\Z,\;m)$ and those with
$(\Delta,n)=(\Z/m\Z,\;2m)$ correspond bijectively. \F
\end{lemma}

\begin{proof}
Only the anchor exchange needs an argument. An $n=2m$ object at odd $m$ has (by
the proof of Theorem~\ref{thm:rigidity}) angles in a coset
$\frac{e}{2m}+\frac1m\Z/\Z$ with $e$ odd. The shift by
$\tfrac12=\frac{m}{2m}$ moves the offset to $\frac{e+m}{2m}$; since $e$ and
$m$ are both odd, $e+m$ is even, so
$\frac{e+m}{2m}=\frac{(e+m)/2}{m}\in\frac1m\Z/\Z$: the image has $n=m$
(angles in $\frac1m\Z/\Z$ with differences generating it). Conversely the
shift moves $\frac1m\Z/\Z$ into the odd coset. (The one-line arithmetic is
nonetheless corroborated on the engine for $m\in\{3,5,7,9\}$ --- ledger
entry~L --- per the project convention that every load-bearing step, however
small, gets a machine check; $T(x^3+2)=x^3-2$ is ledger entry~B.)
\end{proof}

\begin{definition}
For $m\ge1$ let
$\mu_{\mathrm{rel}}(m)=\inf\bigl\{\Mah(O)\;:\;O\text{ relationally
admissible},\ \Delta(O)\cong\Z/m\Z,\ \Mah(O)>1\bigr\}$.
\end{definition}

\begin{theorem}[Reference-free parity floor]\label{thm:parityrel}
\[
\mu_{\mathrm{rel}}(m)=
\begin{cases}
\varphi, & m\text{ even (attained, by }q_{m/2}\text{)},\\[2pt]
\mu(m), & m\text{ odd,}
\end{cases}
\]
where $\mu$ is the absolute floor of~\cite[Thm.~5.1]{CMC}. Statuses: the even
attainment is constructive \F; the identification of the even infimum as
$\varphi$ (rather than a value in $[\theta_0,\varphi)$) is conditional on
\textup{(EG)} (Remark~\ref{rem:gapstatus}). For odd
$m$ the identity
$\mu_{\mathrm{rel}}(m)=\mu(m)$ is \F, so the value inherits~\cite{CMC}:
$\mu_{\mathrm{rel}}(3)=2$ \F, and $\mu_{\mathrm{rel}}(m)=2$ for general odd $m$
is \C{} (to the searched ranges of~\cite[\S6]{CMC}) / \Pl.
\end{theorem}

\begin{proof}
\emph{Even $m=2k$.} Lemma~\ref{lem:qk} gives an object with $\Delta=\Z/m\Z$ and
$\Mah=\varphi$; rigidity makes every competitor charge-admissible, so the
emission gap bounds $\mu_{\mathrm{rel}}(m)\ge\varphi$, with the status
recorded in Remark~\ref{rem:gapstatus}.

\emph{Odd $m$.} By Theorem~\ref{thm:rigidity} the domain splits into the anchor
sectors $n=m$ and $n=2m$. Lemma~\ref{lem:oddfull} identifies the first with the
full family of charge-admissible objects of absolute group $\Z/m\Z$ --- exactly
the domain of $\mu(m)$ --- and Lemma~\ref{lem:twist} maps the second onto the
first bijectively and $\Mah$-preservingly. Hence the infimum over the union
equals $\mu(m)$.
\end{proof}

\begin{remark}[Robustness if \textup{(EG)} fails]\label{rem:robust}
The hypothesis \textup{(EG)} enters exactly once: as the lower bound in the
evaluation of the floors. If \textup{(EG)} were to fail, every statement of
\S\S\ref{sec:relational}--\ref{sec:descent} and
\S\S\ref{sec:inert}--\ref{sec:locking} is unaffected, and within
Theorem~\ref{thm:parityrel} the identity
$\mu_{\mathrm{rel}}(\text{odd})=\mu(\text{odd})$, the even attainment at
$\varphi$ (by $q_k$), and the odd attainment at $2$ (by $x^m-2$) all stand;
only the numerical evaluations of the infima could move below the stated
values. The conditional content of this note is confined to two numbers, not
to any structure.
\end{remark}

\begin{corollary}[The parity bit is intrinsic]\label{cor:paritybit}
The charge--measure coupling of~\cite[Thm.~5.1, Cor.~6.3]{CMC} is a relational
statement: the universal floor $\varphi$ is attainable at nontrivial relational
charge iff the order-$2$ relation of the golden pair
(Lemma~\ref{lem:golden}) embeds in $\Delta(O)$, i.e.\ iff $\tfrac12\in\Delta(O)$,
i.e.\ iff $|\Delta(O)|$ is even. The single parity bit that governs the coupling
is the order of a \emph{relational} class, with no reference ray anywhere; and by
the twist reduction the relational refinement adds no new burden to the
companion's open odd-floor problem. Mechanism \F; universality inherits the
caveats of~\cite[Thm.~5.1, Thm.~6.2]{CMC}.
\end{corollary}

\section{Coherence types on the inadmissible sector}\label{sec:inert}

By Corollary~\ref{cor:nonew} the relational theory adds nothing on admissible
objects. Its genuine domain is the sector the absolute theory discards ---
paradigmatically the Salem numbers, which by~\cite[Lem.~6.1]{CMC} carry no
absolute charge at all, yet by Definition~\ref{def:type} still carry a finite
coherence type. This section makes the type computable and computes it for the
two keystone occupants of the gap $(\varphi,2)$.

\subsection{Shells and the ratio object}

\begin{definition}[Shells]
The \emph{shells} of $O$ are the equivalence classes of conjugates under equal
modulus. A ratio $\alpha/\beta$ is unimodular iff $\alpha,\beta$ share a shell.
\end{definition}

\begin{definition}[Ratio object]\label{def:rat}
For monic $p\in\Z[x]$ of degree $n$ with $p(0)\neq0$, the \emph{ratio
polynomial} is the primitive part of
\[
\Rat_p(x)\;=\;\Res_y\bigl(p(y),\,p(xy)\bigr)\;\in\;\Z[x].
\]
Because $p$ is monic,
\[
\Res_y\bigl(p(y),\,p(xy)\bigr)=\prod_{i}p(x\alpha_i)
=\Bigl(\prod_i\alpha_i\Bigr)^{\!n}\prod_{i,j}\bigl(x-\alpha_j/\alpha_i\bigr),
\qquad \prod_i\alpha_i=(-1)^n p(0)\neq0,
\]
so the degree is exactly $n^2$ (no drop: the leading coefficient is
$(\pm p(0))^n$) and the root multiset is all $n^2$ ordered ratios
$\alpha_j/\alpha_i$. When $p(0)=\pm1$ the companion construction
$\det\bigl(xI-C_p\otimes C_p^{-1}\bigr)$ (an integer matrix) has the same root
multiset and serves as an independent second route (ledger entry~G).
\end{definition}

\medskip
\noindent\textbf{Standing assumption (this section): every $p$ considered
is squarefree --- here each is irreducible.} Lemma~\ref{lem:diag} uses this.

\begin{lemma}[Contact criterion]\label{lem:contact}
Let $\alpha,\beta$ lie in a common shell. Then
$t(\alpha)-t(\beta)\in\tors$ $\iff$ $\alpha/\beta$ is a root of unity $\iff$
some cyclotomic $\Phi_M$ with $\alpha/\beta$ among its roots divides $\Rat_p$.
\F
\end{lemma}

\begin{proof}
$\alpha/\beta$ is unimodular; a unimodular complex number with rational angle
$a/b$ equals $e^{2\pi ia/b}$ and satisfies $z^b=1$. Conversely roots of unity
have rational angle. Divisibility: $\alpha/\beta$ is a root of $\Rat_p$, and if
it is a primitive $M$-th root of unity its minimal polynomial $\Phi_M$ divides
$\Rat_p$ over $\Q$ (Gauss).
\end{proof}

\begin{lemma}[Completeness of the contact scan]\label{lem:complete}
Let $P\in\Z[x]$ be nonzero. If $\Phi_M\mid P$ then $\phi(M)\le\deg P$; and
for every $M\ge1$, $\phi(M)\ge\sqrt{M/2}$. Hence scanning all
$M\le2(\deg P)^2$ decides, completely and in integer arithmetic, which
cyclotomics divide $P$. For $P=\Rat_p$, of degree $n^2$, the bound is
$2n^4$. \F
\end{lemma}

\begin{proof}
The degree bound is $\deg\Phi_M=\phi(M)\le\deg P$. For the totient
bound --- the claim, unambiguously, is $\phi(M)\ge\sqrt{M/2}$, i.e.\
$2\,\phi(M)^2\ge M$ --- write $M=2^aM'$ with $M'$ odd and use
multiplicativity. On the $2$-part, for $a\ge1$,
$\phi(2^a)=2^{a-1}$ while $\sqrt{2^a/2}=2^{(a-1)/2}$; since
$a-1\ge(a-1)/2$ for $a\ge1$ (equality only at $a=1$),
$\phi(2^a)\ge\sqrt{2^a/2}$.
On each odd prime power, $\phi(p^b)=p^{b-1}(p-1)\ge p^{b-1}\sqrt p\ge
p^{b/2}=\sqrt{p^b}$, using $p-1\ge\sqrt p$ for $p\ge3$; multiplying over
the odd part, $\phi(M')\ge\sqrt{M'}$. Assembling: if $a\ge1$,
$\phi(M)=\phi(2^a)\,\phi(M')\ge\sqrt{2^a/2}\cdot\sqrt{M'}=\sqrt{M/2}$;
if $a=0$, $\phi(M)\ge\sqrt M\ge\sqrt{M/2}$. The bound is tight at $M=2$
($\phi(2)=1=\sqrt{2/2}$), so the constant $\tfrac12$ cannot be removed.
Hence $\phi(M)\le\deg P$ forces $M\le2(\deg P)^2$. (The inequality
$2\phi(M)^2\ge M$ is additionally corroborated by exhaustive integer sieve
for all $M\le2\cdot10^{5}$, comfortably past every scan bound used here,
with $M=2$ the unique tight case: ledger entry~N.)
\end{proof}

\begin{lemma}[Diagonal normalization]\label{lem:diag}
If $p$ is squarefree, the multiplicity of $\Phi_1=x-1$ in $\Rat_p$ is exactly
$n$ (the diagonal ratios $\alpha_i/\alpha_i$). Any excess multiplicity, or any
contact with $\Phi_M$, $M\ge2$, certifies a nontrivial torsion ratio. \F
\end{lemma}

\begin{proof}
The diagonal ratios $\alpha_i/\alpha_i=1$ account for $n$ roots of $x-1$;
squarefreeness makes the $\alpha_i$ distinct, so $\alpha_j/\alpha_i=1$ iff
$i=j$, and the multiplicity is exactly $n$.
\end{proof}

\begin{remark}[Scope of the probe]\label{rem:scope}
Cyclotomic contact decides coherence \emph{within shells}. Cross-shell
coherence is equivalent to $\nu=\mu/\overline{\mu}$ being a root of unity for
$\mu=\alpha/\beta$ (note $t(\nu)=2\,t_{\mathrm{rel}}(\alpha,\beta)$, so the
criterion decides rationality exactly while reading the order only up to its
$2$-part); this is again a torsion question about an algebraic number and is
decidable, but it is not needed below: for the ledger's admissible objects the
angles are available in closed form and the full type is settled symbolically,
while for a Salem number the shells are $\{\tau\}$, $\{1/\tau\}$ and the unit
circle, so by Corollary~\ref{cor:salemgeom} the contact scan is a complete
decision for the only undetermined block. The shell signature is a strictly
coarser invariant than the full type: $x^4+5x^2+5$ and the catalog seed
$K=x^4+5x^2-5$ have identical contact signatures $\{\Phi_1^4,\Phi_2^4\}$
(ledger entries~E, F) while their full relational groups are $\Z/2\Z$ and
$\Z/4\Z$; and the probe is gauge-blind by construction --- $x^3-2$ (absolute
$\Z/3\Z$) and $x^3+2$ (absolute $\Z/6\Z$) return identical signatures
$\{\Phi_1^3,\Phi_3^3\}$, both reading the common relational group $\Z/3\Z$
(ledger entries~A, B). \F
\end{remark}

\subsection{The two keystone occupants are relationally inert}

\begin{definition}[Relationally inert]\label{def:inert}
A conjugation-closed object is \emph{relationally inert} if every coherence
class is either contained in the real axis or a singleton --- equivalently,
if the only coherent pairs of distinct conjugates are pairs of real
conjugates. (Real conjugates are always mutually coherent, their angles
lying in $\{0,\tfrac12\}$; inertness says no other coherence exists.)
\end{definition}

\begin{theorem}[Inertness of $\beta_4$ and Lehmer's number]\label{thm:inert}
Let $\beta_4=x^4-x^3-x^2-x+1$ (the minimal degree-$4$ Salem number,
$\Mah=1.72208\ldots$) and let $L$ be Lehmer's polynomial
$x^{10}+x^9-x^7-x^6-x^5-x^4-x^3+x+1$ ($\Mah=1.17628\ldots$). The complete
contact scans return
\[
\Rat_{\beta_4}:\ \{\Phi_1^{\,4}\}\qquad\text{and}\qquad
\Rat_{L}:\ \{\Phi_1^{\,10}\}
\]
and nothing else (complete scan bounds $M\le512=2\cdot16^2$ and
$M\le20000=2\cdot100^2$: Lemma~\ref{lem:complete} applied to
$\deg\Rat_{\beta_4}=16$ and $\deg\Rat_L=100$). Combined
with Corollary~\ref{cor:salemgeom}, the coherence types are: rational block
$\{\tau,1/\tau\}$ with trivial class group, plus $2s$ mutually incoherent
circle singletons ($s=1$ for $\beta_4$; $s=4$ for $L$). In words: \emph{both
objects are relationally inert (Definition~\ref{def:inert}); the only rational
phase relations they possess are the tautological diagonal and the
positivity-forced real pair $(\tau,1/\tau)$.}
\F{} per instance (ledger entries~G, H)
\end{theorem}

\begin{theorem}[Extension across degrees $4$--$10$]\label{thm:inertmore}
Let $S_6=x^6-x^4-x^3-x^2+1$ and $S_8=x^8-x^5-x^4-x^3+1$. All four polynomials
$\beta_4$, $S_6$, $S_8$, $L$ are Salem, certified exactly: each is reciprocal
and irreducible, and its trace polynomial --- the $T\in\Z[z]$ with
$p(x)=x^{\deg p/2}\,T(x+1/x)$ --- has, by Sturm count, exactly one real root
in $(2,\infty)$, none in $(-\infty,-2]$, and the remaining $\deg p/2-1$ in
$(-2,2)$ (ledger entry~O). The complete contact scans of $\Rat_{S_6}$ and
$\Rat_{S_8}$ return $\{\Phi_1^{\,6}\}$ and $\{\Phi_1^{\,8}\}$ respectively
(degrees $36$ and $64$; bounds $2592$ and $8192$), so both are relationally
inert (Definition~\ref{def:inert}; ledger entry~P). \F{} per instance
\end{theorem}

\begin{theorem}[Exhaustive degree-$12$ census]\label{thm:census}
Consider all $3^6=729$ monic reciprocal polynomials
$x^{12}+c_1x^{11}+\dots+c_6x^6+\dots+c_1x+1$ with
$c_1,\dots,c_6\in\{-1,0,1\}$. The sign twist $p(x)\mapsto p(-x)$ preserves
the family and fixes $\Rat_p$ (a global sign cancels in every ratio), so the
family splits into twist-classes counted by Burnside's lemma: the involution
fixes exactly the $3^3=27$ coefficient vectors with vanishing odd
coefficients, giving $(729+27)/2=378$ orbits --- an internal count check
reproduced by the enumeration. Exactly $37$ twist-classes are
Salem, each certified as in Theorem~\ref{thm:inertmore}; the complete
contact scan of every one ($\deg\Rat_p=144$, bound $41472=2\cdot144^2$)
returns $\{\Phi_1^{\,12}\}$: \emph{every certified Salem number in the
family is relationally inert} (ledger entry~T). \F{} per instance; the
family statement is a complete finite verification, \C
\end{theorem}

\begin{remark}[Scope of the census]\label{rem:censusscope}
The $\{-1,0,1\}$ family was chosen for exhaustive tractability, not
representativeness; with Corollary~\ref{cor:allinert} below in hand the
census is corroboration rather than evidence. Five further degree-$12$
Salem instances with a coefficient outside $\{-1,0,1\}$ were spot-checked
and are likewise inert (ledger entry~V).
\end{remark}

\begin{corollary}[Sharpened exclusion]\label{cor:sharpexcl}
For these two instances, \cite[Lem.~6.1]{CMC} strengthens: not only does no
finite absolute charge group exist --- no nontrivial rational phase relation
exists at all among the conjugates off the real axis. The would-be occupants of
the odd-charge gap $(\varphi,2)$ studied in~\cite[Thm.~6.2]{CMC} are, at least
at the two keystones, maximally far from carrying phase structure in either the
absolute or the relational sense. Inertness is a positive finding, not the
absence of one: each scan excludes \emph{every} candidate torsion relation
among the $n^2$ ratios --- an infinite family of cyclotomic hypotheses --- in
a single finite certificate, and it is exactly this that upgrades
\cite[Lem.~6.1]{CMC} from ``no finite charge group'' to ``no rational phase
structure whatsoever off the real axis'' for these witnesses. Every Salem
number is in fact relationally inert (Corollary~\ref{cor:allinert} below);
the per-instance decidability of
Lemmas~\ref{lem:contact}--\ref{lem:complete} remains the general
instrument beyond the pinned-modulus case.
\end{corollary}

\subsection{A general theorem: modulus pinning}

\begin{theorem}[Modulus pinning]\label{thm:pinning}
Let $p\in\Z[x]$ be irreducible with a root $\alpha_0$ whose modulus is
attained by no other root of $p$. Then no two distinct roots of $p$ differ
by a root-of-unity factor. More generally, if $q\in\Z[x]$ is irreducible,
$q\neq p$, and no root of $q$ has modulus $|\alpha_0|$, then no root of $p$
and root of $q$ differ by a root-of-unity factor. \F
\end{theorem}

\begin{proof}
Suppose $\alpha\neq\beta$ are roots of $p$ with $\alpha/\beta=\zeta$ a root
of unity. Let $G$ be the Galois group of the splitting field of $p$; by
irreducibility $G$ acts transitively on the roots, so some $\sigma\in G$
has $\sigma\alpha=\alpha_0$. Then
$\sigma\alpha/\sigma\beta=\sigma(\zeta)$ is again a root of unity, so
$|\sigma\beta|=|\alpha_0|$; since $\sigma\beta$ is a root of $p$ and
$\alpha_0$ is the only root of that modulus, $\sigma\beta=\alpha_0
=\sigma\alpha$, whence $\alpha=\beta$ --- a contradiction. For the mixed
statement, work in a Galois extension containing both splitting fields: the
restriction map onto $\mathrm{Gal}$ of the splitting field of $p$ is
surjective, so again some $\sigma$ carries the $p$-root of the pair to
$\alpha_0$, and the image of the $q$-root would be a root of $q$ of modulus
$|\alpha_0|$, which does not exist.
\end{proof}

\begin{corollary}[Every Salem number is relationally inert]
\label{cor:allinert}
Every Salem number is relationally inert; open problem P1 of the earlier
revisions of this note is resolved in the negative, and
Theorems~\ref{thm:inert}, \ref{thm:inertmore}, \ref{thm:census} become
corroborations of a general fact. \F
\end{corollary}

\begin{proof}
The modulus $\tau$ is attained only by $\tau$, so
Theorem~\ref{thm:pinning} applies. A non-singleton class not contained in
the real axis would contain either two distinct circle conjugates ---
coherent iff their ratio is a root of unity (Lemma~\ref{lem:contact}),
excluded --- or a real and a circle conjugate $\lambda$: coherence would
force $t(\lambda)\in\Q/\Z$, making $\lambda/\overline{\lambda}$
(unimodular, rational angle) a root of unity, again excluded. So every
class is real-contained or a singleton, with real block $\{\tau,1/\tau\}$
of trivial class group.
\end{proof}

\begin{corollary}[No two Salem numbers circle-lock]\label{cor:nolock}
No two distinct Salem numbers are circle-locked; indeed no root of one and
root of another differ by a root-of-unity factor. Open problem P2 is
resolved in the negative; Theorem~\ref{thm:notlocked} and
Remark~\ref{rem:allpairs} are corroborations. \F
\end{corollary}

\begin{proof}
Apply the mixed statement with $\alpha_0=\tau_p$: a root of $q$ of modulus
$\tau_p$ would force $\tau_p\in\{\tau_q,1/\tau_q,1\}$, i.e.\
$\tau_p=\tau_q$, hence the same minimal polynomial, $p=q$.
\end{proof}

\begin{corollary}[Pisot numbers]\label{cor:pisot}
No two distinct conjugates of a Pisot number differ by a root-of-unity
factor, and no non-real conjugate is coherent with a real one. (The full
type of a Pisot number can still carry mirrored cross-shell classes;
deciding those is the $\nu$-criterion of Remark~\ref{rem:scope}.) \F
\end{corollary}

\begin{proof}
The dominant root pins its modulus; the real--non-real piece follows as in
Corollary~\ref{cor:allinert}: coherence with a real conjugate would make
$t(\lambda)$ rational and $\lambda/\overline{\lambda}$ a torsion ratio.
\end{proof}

\begin{example}[Sharpness, and an irreducible non-inert witness]
\label{ex:twistshell}
The pinning hypothesis is necessary: $x^4-2$, with all four roots on one
shell, is irreducible with torsion ratios $\zeta_4^{\,j}$ (ledger
entry~C). And irreducible \emph{inadmissible} non-inert objects exist: for
$p(x)=q(x^k)$ with $q$ irreducible carrying a non-real root at irrational
angle, each fiber of $k$-th roots is a size-$k$ coherence class with class
group $\Z/k\Z$ at an irrational offset. Concretely,
$p=x^4+x^2+2=q(x^2)$ with $q=x^2+x+2$: $p$ is irreducible; $q$'s angle is
irrational (complete scan of $\Rat_q$: $\{\Phi_1^{\,2}\}$, ledger
entry~U), so $p$ is inadmissible; and $\Rat_p$ has signature
$\{\Phi_1^{\,4},\Phi_2^{\,4}\}$ --- the fibers $\pm\sqrt{\rho}$,
$\pm\sqrt{\overline\rho}$ are non-real size-$2$ classes, so $p$ is not
inert; $\Mah(p)=2$. Salem polynomials can never arise this way: $q(x^k)$
with $k\ge2$ would duplicate the maximal modulus --- pinning and the
twisted-shell mechanism exclude each other. \F
\end{example}

\begin{remark}[Degeneracy, and the literature]\label{rem:degen}
In the language of linear recurrence sequences, distinct characteristic
roots with a root-of-unity ratio are precisely the \emph{degenerate} case
(see~\cite{EPSW}); Theorem~\ref{thm:pinning} says that irreducible
polynomials with a pinned modulus --- in particular all Salem and Pisot
polynomials --- are non-degenerate. The argument is elementary and may well
be classical in that literature; we have not located a statement covering
the mixed two-polynomial case or the coherence-type packaging, and we
record the proof self-containedly. What the relational framework adds
beyond non-degeneracy is the rest of the type: offsets, mirrored classes,
the real and mixed bookkeeping, the descent laws, and the shell-complete
decision procedure that corroborated the theorem on forty-six instances
before it was proved.
\end{remark}

\section{Cross-object locking}\label{sec:locking}

Rigidity kills within-object novelty; it says nothing about relations
\emph{between} different oscillator systems. By
Proposition~\ref{prop:locking}, two objects can each be relationally inert yet
be locked --- a charge that lives purely in the relationship, with no per-object
counterpart. For Salem numbers the interesting case is circle-locking, and it is
decidable:

\begin{definition}[Mixed ratio object]
For monic $p,q\in\Z[x]$ with nonzero constant terms, the primitive part of
$\Rat_{p,q}(x)=\Res_y\bigl(q(y),\,p(xy)\bigr)$ has root multiset
$\{\alpha/\beta:\ p(\alpha)=0,\ q(\beta)=0\}$, of degree $(\deg p)(\deg q)$.
\end{definition}

For two Salem numbers the unimodular cross-ratios are exactly the
circle-by-circle ones (the real conjugates of distinct Salem numbers never share
a modulus: $\tau_p=\tau_q$ is excluded by distinctness and
$\tau_p=1/\tau_q$ by $\tau>1$), so the contact scan on $\Rat_{p,q}$ is a
complete decision for circle-locking; contact with $\Phi_1$ would mean a shared
root, excluded when $\gcd(p,q)=1$. (Corollary~\ref{cor:nolock} now settles
circle-locking for the entire Salem class; the scans below stand as the
original exact corroboration.)

\begin{theorem}[$\beta_4$ and Lehmer's number are not circle-locked]
\label{thm:notlocked}
$\gcd(\beta_4,L)=1$, and the complete contact scan of the degree-$40$ mixed
ratio object $\Rat_{L,\beta_4}$ (complete bound $M\le3200=2\cdot40^2$:
Lemma~\ref{lem:complete} applied to $\deg\Rat_{L,\beta_4}=40$) returns
\emph{no} cyclotomic
contact whatsoever. Hence no circle conjugate of Lehmer's number differs from a
circle conjugate of $\beta_4$ by a rational angle: the two keystone occupants of
$(\varphi,2)$ are mutually as incoherent as each is internally. \F{} per
instance (ledger entry~I)
\end{theorem}

\begin{remark}
The trivial real locking $\tau_{\beta_4}\approx\tau_L$ (both at angle $0$) is of
course present; it is the rational-block locking that
Proposition~\ref{prop:locking} discounts. What Theorem~\ref{thm:notlocked} rules
out is the informative kind.
\end{remark}

\begin{remark}[All pairs among the certified instances]\label{rem:allpairs}
The same verdict holds for every pair among the four instances of
Theorem~\ref{thm:inertmore}: the five remaining mixed ratio objects
($\beta_4\times S_6$, $\beta_4\times S_8$, $S_6\times S_8$, $L\times S_6$,
$L\times S_8$; degrees $24,32,48,60,80$, complete bounds $1152$--$12800$)
carry no cyclotomic contact (ledger entry~Q). No two of the four certified
Salem numbers are circle-locked. \F{} per instance
\end{remark}

\begin{example}[A non-inert object, computed: $\beta_4\otimes\beta_4$]
\label{ex:tensorsq}
Write $\beta_4$'s conjugates as $\tau,1/\tau,\lambda,\overline\lambda$, with
$\lambda$ on the unit circle at irrational angle $\theta$. The $16$
eigenvalues of $\beta_4\otimes\beta_4$ sort into the coherence type: a
rational block of six positive reals $\{1,1,1,1,\tau^2,\tau^{-2}\}$ at
offset $0$; a mirrored pair of size-$4$ classes at offsets $\pm\theta$,
namely $\{\tau\lambda,\tau\lambda,\tau^{-1}\lambda,\tau^{-1}\lambda\}$ and
its conjugate --- two distinct values, each of multiplicity two, sharing a
single angle, so the class group is trivial while the class is non-real of
size $4$; and a mirrored singleton pair
$\{\lambda^2\},\{\overline\lambda^{\,2}\}$ at offsets $\pm2\theta$. By
Definition~\ref{def:inert} the object is \emph{not} inert: the
offset-cancellation mechanism of Proposition~\ref{prop:tensor}(iii)
manufactures both a rational block and non-real coherence from an inert
factor. Engine corroboration (ledger entry~S): the characteristic polynomial
$F$ of the $16\times16$ Kronecker square $C_{\beta_4}\otimes C_{\beta_4}$
has $(x-1)$-multiplicity exactly $4$, $\deg\gcd(F,F')=7$ (exactly the
multiplicity pattern $1^4\cdot(\text{four doubles})$), and exactly $3$
distinct real roots, all positive --- Sturm counts, no floats. \F
\end{example}

\section{Methodology}\label{sec:method}

\textbf{Exact arithmetic, strengthened.} The companion note decided angle
questions at $25$--$35$ decimal digits with a tolerance discipline. The
relational questions of this note are \emph{torsion} questions, and the entire
verification path here is integer and rational polynomial arithmetic: resultants
over $\Z[x]$, primitive parts, and exact trial division by cyclotomic
polynomials. No floating-point number appears anywhere --- not in a decision,
not in a displayed value. The completeness certificate of
Lemma~\ref{lem:complete} replaces any notion of numerical tolerance.

\textbf{Engine.} \texttt{sympy}~1.14.0 on Python~3.12.3. The consolidation
ledger of Appendix~A was produced by the scripts
\texttt{relational\_charge\_probe.py} and \texttt{supplement\_for\_paper.py}
(Appendix~B lists the probe core), each rerun in full immediately prior to
writing; every numerical claim in the text is grounded in a ledger entry. SHA-256
digests of the deliverables (source, PDF, scripts) accompany the artifacts
in \texttt{SHA256SUMS}.

\textbf{Second engine.} Following review, the complete contact signatures
of ledger entries A--I, together with rows P, Q, U, and V, were
recomputed independently in PARI/GP~2.15.4 via
\texttt{polresultant}, integer factorization, and \texttt{poliscyclo}; all
twenty-three signatures agree with the \texttt{sympy} computations (ledger
entry~M). The two stacks share no computer-algebra code: \texttt{sympy}'s
resultants and factorizations are pure Python, PARI's are its own C library.

\textbf{Two-route corroboration.} For $\beta_4$ the ratio object was computed
independently by the resultant of Definition~\ref{def:rat} and as the
characteristic polynomial of the $16\times16$ integer matrix
$C_{\beta_4}\otimes C_{\beta_4}^{-1}$; the results agree (ledger entry~G).

\textbf{A harness incident, recorded.} The first assertion written for the
control $q_2$ predicted $\Phi_2$-multiplicity $2$ and failed; the engine was
right and the assertion wrong --- ratios are \emph{ordered} pairs, so each
antipodal pair contributes twice. Per the project discipline that load-bearing
corrections are pinned rather than papered over, the incident is recorded here
and the corrected expectation ($\Phi_2^{\,4}$) is ledger entry~D.

\textbf{Limits of self-verification.} The cross-engine run was performed by
the same operator as the primary computation; the stacks are disjoint, the
operator is not. Independent replication by external researchers is welcomed:
the scripts, inputs, expected signatures, and digests are provided for
exactly this purpose (\texttt{relational\_charge\_probe.py},
\texttt{supplement\_for\_paper.py}, \texttt{supplement\_round3.py},
\texttt{census\_deg12.py}, \texttt{supplement\_round6.py},
\texttt{cross\_check.gp}, \texttt{cross\_round6.gp},
\texttt{SHA256SUMS}); the project repository is DOI-archived at
\url{10.5281/zenodo.21121863}; the scripts above are distributed with this
manuscript and intended for the deposit's next version, enabling
verification by a distinct operator. Concretely (expected
outputs are the signatures of Appendix~A; any deviation is a finding):
\begin{itemize}
\item \texttt{python3 relational\_charge\_probe.py} --- entries A--H under
assertions (e.g.\ $\{\Phi_1^{\,10}\}$ for Lehmer; about one second);
\item \texttt{python3 supplement\_round3.py} --- entries O--S;
\item \texttt{python3 supplement\_round6.py} --- entries U--V;
\item \texttt{python3 census\_deg12.py} --- the census of
Theorem~\ref{thm:census} (about two minutes);
\item \texttt{gp -q cross\_check.gp} and \texttt{gp -q cross\_round6.gp}
--- the twenty-three cross-engine signatures.
\end{itemize}

\section{Scope, limitations, and open problems}\label{sec:open}

\begin{center}
\small
\begin{tabular}{p{4.9cm}p{5.6cm}p{3.6cm}}
\hline
Result & Statement & Status\\
\hline
Cyclicity (Prop.~\ref{prop:cyclic}) & class/relational groups cyclic & \F\\
Gauge (Prop.~\ref{prop:gauge}) & absolute $=$ relational $+$ ray & \F\\
Rigidity (Thm.~\ref{thm:rigidity}) & relational adm.\ $\Rightarrow$ absolute; $n\in\{m,2m\}$ & \F\\
Structure (Thm.~\ref{thm:structure}) & rational block $+$ mirrored offset classes & \F\\
Odd anchors full (Lem.~\ref{lem:oddfull}) & $n$ odd $\Rightarrow\Delta=\Z/n\Z$ & \F\\
Descent (\S\ref{sec:descent}) & $\lcm$ law, CRT, safe $\Leftrightarrow$ coprime & \F\\
Internalization (Prop.~\ref{prop:tensor}iii) & $\otimes$ creates rational blocks by offset cancellation & \F\\
Golden relation (Lem.~\ref{lem:golden}) & $t_{\mathrm{rel}}(\varphi,\varphi')=\tfrac12$ & \F\\
Twist (Lem.~\ref{lem:twist}) & anchor sectors exchanged at odd $m$, $\Mah$ fixed & \F\\
Parity floor (Thm.~\ref{thm:parityrel}) & $\mu_{\mathrm{rel}}(\text{even})=\varphi$ attained & \F\\
 & $\mu_{\mathrm{rel}}(\text{odd})=\mu(\text{odd})$; value $2$ & \F{} at $3$; \C/\Pl{} general\\
 & lower bound $\varphi$ & conditional on (EG), Rem.~\ref{rem:gapstatus}\\
Inertness (Thm.~\ref{thm:inert}) & $\beta_4$, Lehmer relationally inert & \F{} per instance\\
Extension (Thm.~\ref{thm:inertmore}) & $S_6$, $S_8$ certified Salem; inert & \F{} per instance\\
Pinning (Thm.~\ref{thm:pinning}) & pinned modulus forbids torsion ratios & \F\\
Salem inertness (Cor.~\ref{cor:allinert}) & every Salem number inert; P1 closed & \F\\
No locking (Cor.~\ref{cor:nolock}) & no two Salem numbers circle-lock; P2 closed & \F\\
Not locked (Thm.~\ref{thm:notlocked}, Rem.~\ref{rem:allpairs}) & six pairwise scans & corroboration\\
Census (Thm.~\ref{thm:census}) & $37/37$ degree-$12$ family Salems inert & corroboration; \C{} family\\
Cross-engine (App.~A, entry M) & PARI/GP reproduces signatures A--I, P, Q & \C\\
\hline
\end{tabular}
\end{center}

\medskip
\noindent\textbf{Open problems.}
\begin{enumerate}
\item[(P1)] \textbf{Internal Salem coherence --- resolved in the negative.}
Corollary~\ref{cor:allinert}: every Salem number is relationally inert.
The forty-six certified instance decisions of this note --- the four
canonical test cases of degrees $4$--$10$, the exhaustive degree-$12$
family, and the five larger-coefficient spot checks (ledger entries G, H,
P, T, V) --- stand as exact corroboration, obtained \emph{before} the
theorem. The successor problem is the characterization: \textbf{(P1$'$)}
which irreducible integer polynomials carry nontrivial coherence? A torsion
ratio $\zeta_M$ forces $\Q(\zeta_M)$ into the splitting field; a pinned
modulus forbids torsion (Theorem~\ref{thm:pinning}); twisted shells
$q(x^k)$ suffice (Example~\ref{ex:twistshell}); and composite objects
manufacture coherence freely (Example~\ref{ex:tensorsq}). In recurrence
language: characterize degenerate irreducibles beyond the
pinned/twisted-shell dichotomy (cf.~\cite{EPSW}), and, beyond torsion,
classify the cross-shell mirrored classes the $\nu$-criterion governs. \Op
\item[(P2)] \textbf{Cross-Salem locking --- resolved in the negative.}
Corollary~\ref{cor:nolock}: the circle angles of distinct Salem numbers,
viewed in $(\R/\Z)/\Q$, never collide. Theorem~\ref{thm:notlocked} and
Remark~\ref{rem:allpairs} (six pairs) are exact corroboration. \F
\item[(P3)] \textbf{The odd floor.} Prove $\mu_{\mathrm{rel}}(m)=2$ for all odd
$m$. By Lemma~\ref{lem:twist} this is \emph{exactly} the companion's principal
open problem~\cite[\S8]{CMC} --- the relational refinement adds no new burden,
since the new anchor sector reduces by the twist. \Op
\item[(P4)] \textbf{Is parity the whole coupling?} Does
$\mu_{\mathrm{rel}}$ depend on $\Delta$ through anything besides the single bit
$2\mid m$? The values known and conjectured (Theorem~\ref{thm:parityrel})
suggest not. \Op
\end{enumerate}

\medskip
\noindent\textbf{Caveat on the larger programme.} As in~\cite{CMC}, nothing here
touches Lehmer's problem itself or the substrate's deeper conjectures. What this
note establishes is internal to the phase structure: the absolute charge was a
gauge-fixed presentation of a relational invariant; integrality is what fixes
the gauge; the parity coupling survives the deletion of the reference ray
unchanged; and on the sector the absolute theory discards, the relational
theory is nontrivial, finite, decidable --- and, at the two keystones of the gap
$(\varphi,2)$, maximally empty.

\appendix

\section{Consolidation ledger}\label{app:ledger}

The following were regenerated in exact arithmetic immediately prior to writing
(engine of \S\ref{sec:method}); each grounds a claim in the text. ``Contacts''
means the complete cyclotomic-contact signature of the stated ratio object
(Lemmas~\ref{lem:contact}--\ref{lem:complete}), written
$\{\Phi_{M}^{\,\mathrm{mult}}\}$. Backing scripts:\\
\centerline{\texttt{relational\_charge\_probe.py}\quad and\quad
\texttt{supplement\_for\_paper.py}.}

\begin{center}
\small
\begin{tabular}{clp{9.2cm}}
\hline
 & object & exact result\\
\hline
A & $x^3-2$ & contacts $\{\Phi_1^{\,3},\Phi_3^{\,3}\}$; $\Delta=\Z/3\Z$, anchor
$n=3$, $\Mah=2$.\\[2pt]
B & $x^3+2$ & contacts $\{\Phi_1^{\,3},\Phi_3^{\,3}\}$, identical to~A: probe is
gauge-blind ($\Delta=\Z/3\Z$; anchor $n=6$); twist $T(x^3+2)=x^3-2$.\\[2pt]
C & $x^4-2$ & contacts $\{\Phi_1^{\,4},\Phi_2^{\,4},\Phi_4^{\,4}\}$: relational
recovery of $\Z/4\Z$ with no reference ray.\\[2pt]
D & $q_2=x^4+x^2-1$ & contacts $\{\Phi_1^{\,4},\Phi_2^{\,4}\}$ (ordered-pair
multiplicities; see \S\ref{sec:method}); full $\Delta=\Z/4\Z$ settled
symbolically; $\Mah=\varphi$.\\[2pt]
E & $x^4+5x^2+5$ & irreducible (Eisenstein at $5$); $x^2$-roots
$\frac{-5\pm\sqrt5}{2}$ both negative, so all roots purely imaginary; anchor
$n=4$, $\Delta=\Z/2\Z$ (group drop), $\Mah=5$; contacts
$\{\Phi_1^{\,4},\Phi_2^{\,4}\}$.\\[2pt]
F & $K=x^4+5x^2-5$ & contacts $\{\Phi_1^{\,4},\Phi_2^{\,4}\}$, identical to~E,
while the full type differs: angles $\{0,\tfrac12,\tfrac14,\tfrac34\}$, so
$\Delta=\Z/4\Z$ --- the shell signature is strictly coarser than the type.\\[2pt]
G & $\beta_4=x^4-x^3-x^2-x+1$ & contacts $\{\Phi_1^{\,4}\}$ only (complete scan
$M\le512$): circle block inert. Two-route: resultant route $=$ characteristic
polynomial of $C_{\beta_4}\otimes C_{\beta_4}^{-1}$.\\[2pt]
H & Lehmer $L$, degree $10$ & $\deg\Rat_L=100$; contacts $\{\Phi_1^{\,10}\}$
only (complete scan $M\le20000$, $1.1$\,s): all $8$ circle conjugates mutually
inert.\\[2pt]
I & $\beta_4\times L$ mixed & $\gcd(\beta_4,L)=1$; $\deg\Rat_{L,\beta_4}=40$;
contacts $=\varnothing$ (complete scan $M\le3200$): not circle-locked.\\[2pt]
J & golden pair & $\varphi/\varphi'=-\varphi^2$: negative real, so
$t_{\mathrm{rel}}(\varphi,\varphi')=\tfrac12$.\\[2pt]
K & $y^2+y-1$ & roots $1/\varphi>0$ and $-\varphi<0$: the real $\pm$ pair
behind Lemma~\ref{lem:qk}, all $k$.\\[2pt]
L & coset arithmetic & $e$ odd, $m$ odd $\Rightarrow(e+m)/2m\in\frac1m\Z/\Z$;
engine corroboration of Lemma~\ref{lem:twist}'s one-line step,
$m\in\{3,5,7,9\}$, all odd $e<2m$.\\[2pt]
M & cross-engine replication & contact signatures of entries A--I, P, Q, U,
and V recomputed in PARI/GP~2.15.4 (\texttt{polresultant} $+$
\texttt{poliscyclo}); all twenty-three identical to the \texttt{sympy}
signatures.\\[2pt]
N & totient bound & $2\,\phi(M)^2\ge M$ verified by exact integer sieve for
all $1\le M\le2\cdot10^{5}$; unique tight case $M=2$
(Lemma~\ref{lem:complete}).\\[2pt]
O & Salem certification & $\beta_4$, $S_6$, $S_8$, $L$: each reciprocal and
irreducible; trace-polynomial Sturm counts $(1,0,\deg p/2-1)$ in
$(2,\infty)$, $(-\infty,-2]$, $(-2,2)$ --- all four Salem, exactly
(Theorem~\ref{thm:inertmore}).\\[2pt]
P & $S_6$, $S_8$ inertness & $\Rat_{S_6}$ (deg $36$, bound $2592$): contacts
$\{\Phi_1^{\,6}\}$; $\Rat_{S_8}$ (deg $64$, bound $8192$): contacts
$\{\Phi_1^{\,8}\}$.\\[2pt]
Q & pairwise locking batch & the five remaining pairs among
$\{\beta_4,S_6,S_8,L\}$: mixed ratio objects of degrees $24,32,48,60,80$
(bounds $1152$--$12800$); contacts $=\varnothing$ in every case.\\[2pt]
R & quadratic identities & $(1+\sqrt5)/2=\varphi$ and
$(3+\sqrt5)/2=\varphi^2$ (Appendix~\ref{app:quadEG}).\\[2pt]
S & $\beta_4\otimes\beta_4$ (non-inert witness) & charpoly $F$ of the
$16\times16$ Kronecker square: $(x-1)$-multiplicity $=4$;
$\deg\gcd(F,F')=7$ (pattern $1^4\cdot$four doubles); $3$ distinct real
roots, all positive (Sturm). Structure forced symbolically; single-engine
run (Example~\ref{ex:tensorsq}).\\[2pt]
T & degree-$12$ census & $729$ reciprocal $\{-1,0,1\}$-polynomials; $378$
twist-classes ($27$ twist-fixed; $378=(729+27)/2$); $37$ certified Salem;
all $37$ complete scans ($\deg144$, bound $41472$) return
$\{\Phi_1^{\,12}\}$. Runtime $93$\,s (\texttt{census\_deg12.py});
cross-engine for the four canonical instances via entry~M.\\[2pt]
U & twisted-shell witness & $q=x^2+x+2$: $\Rat_q$ contacts
$\{\Phi_1^{\,2}\}$ (angle irrational, complete); $p=x^4+x^2+2=q(x^2)$:
irreducible, $\Rat_p$ contacts $\{\Phi_1^{\,4},\Phi_2^{\,4}\}$,
$\Mah=2$ --- non-inert (Example~\ref{ex:twistshell}).\\[2pt]
V & larger-coefficient spot checks & first five degree-$12$ Salem
twist-classes with a coefficient outside $\{-1,0,1\}$ (lexicographic in
$\{-2,\dots,2\}^6$): all five scans return $\{\Phi_1^{\,12}\}$
(Remark~\ref{rem:censusscope}).\\
\hline
\end{tabular}
\end{center}

\section{The probe core}\label{app:probe}

The complete decision procedure of
Lemmas~\ref{lem:contact}--\ref{lem:complete}, as executed. Everything is
integer polynomial arithmetic; \texttt{x, y} are \texttt{sympy} symbols.

\begin{verbatim}
def ratio_poly(p_expr):
    R = sp.resultant(p_expr.subs(x, y), sp.expand(p_expr.subs(x, x*y)), y)
    return sp.Poly(sp.expand(R), x).primitive()[1]

def totients_upto(N):
    phi = list(range(N + 1))
    for i in range(2, N + 1):
        if phi[i] == i:                      # i prime
            for j in range(i, N + 1, i):
                phi[j] -= phi[j] // i
    return phi

def cyclotomic_contacts(P):
    d = P.degree()
    N = 2 * d * d                # complete: Phi_M | P forces M <= 2 d^2
    phi = totients_upto(N)
    hits = {}
    for m in range(1, N + 1):
        if phi[m] > d:
            continue
        cm = sp.Poly(sp.cyclotomic_poly(m, x), x)
        if sp.rem(P, cm).is_zero:  # Phi_m irreducible: divides, or no contact
            mult, T = 0, P
            while True:
                q, r = sp.div(T, cm)
                if r.is_zero:
                    mult, T = mult + 1, q
                else:
                    break
            hits[m] = mult
    return hits
\end{verbatim}

\noindent The PARI/GP cross-check (ledger entries M, P, Q) is the following
independent implementation; its expected outputs are exactly the contact
signatures of Appendix~A.

\begin{verbatim}
contacts(R) = {
  my(F = factor(R), v = List());
  for(i = 1, matsize(F)[1],
    if(type(F[i,1]) == "t_POL" && poldegree(F[i,1]) > 0,
      my(n = poliscyclo(F[i,1]));
      if(n, listput(v, [n, F[i,2]]))));
  vecsort(Vec(v))
}
rat(p)      = polresultant(subst(p, x, 'y), subst(p, x, x*'y), 'y);
mixed(p, q) = polresultant(subst(q, x, 'y), subst(p, x, x*'y), 'y);
\end{verbatim}

\noindent The mixed probe replaces the first argument of the resultant by the
second polynomial: $\Rat_{p,q}=\Res_y\bigl(q(y),\,p(xy)\bigr)$.

\section{Inherited statements from the companion note}\label{app:inherited}

Reference~\cite{CMC} is a technical note of the same project
(the \texttt{math-research-pipelines} corpus), distributed alongside this
one and archived with the repository at DOI
\url{10.5281/zenodo.21121863}. So that the present note can be
read, and its dependence structure audited, without~\cite{CMC} in hand, the
table below lists every
statement imported from it, the status \emph{as tagged there}, and where it
enters this note. Nothing is silently promoted; items this note re-derives
are marked.

\begin{center}
\small
\begin{tabular}{lp{5.4cm}p{3.5cm}p{3.3cm}}
\hline
\cite{CMC} item & statement & status in \cite{CMC} & used here\\
\hline
Def.~2.4 & absolute admissibility; charge group $\Z/n\Z$; charge $\chi_n$ &
definition & Def.~\ref{def:absolute}\\[2pt]
Thm.~3.1 & charge groups are finite cyclic & \F &
Prop.~\ref{prop:cyclic} (re-derived)\\[2pt]
Thm.~3.2 & $x^n-2$ realizes $\Z/n\Z$, $\Mah=2$ & \F &
Thm.~\ref{thm:rigidity}, Thm.~\ref{thm:parityrel}\\[2pt]
Thm.~3.3 & CRT via Adams operators & \F & Prop.~\ref{prop:adams}\\[2pt]
Thm.~4.2 & $\lcm$ law under $\otimes$ & \F &
Prop.~\ref{prop:tensor} (re-derived)\\[2pt]
Thm.~4.3 & safe composition $\Leftrightarrow$ coprime & \F &
Prop.~\ref{prop:tensor}\\[2pt]
Prop.~2.3 & emission gap: $\Mah\in\{1\}\cup[\varphi,\infty)$ for
admissible objects & \F{} (quadratics); \C/\Pl{} (general); named
\textup{(EG)} here &
Rem.~\ref{rem:gapstatus}, Thm.~\ref{thm:parityrel}\\[2pt]
Thm.~5.1 & absolute parity floor $\mu(\mathrm{even})=\varphi$,
$\mu(\mathrm{odd})=2$ & even \F; odd \F{} at $n=3$, \C/\Pl{} general &
Thm.~\ref{thm:parityrel}\\[2pt]
Lem.~5.2 & $\pi$-ray obstruction for the golden seed & \F &
Lem.~\ref{lem:golden}, Cor.~\ref{cor:paritybit}\\[2pt]
Lem.~6.1 & Salem numbers are charge-inadmissible & \F &
Cor.~\ref{cor:salemgeom}, \S\ref{sec:inert}\\[2pt]
Thm.~6.2 & odd gap $(\varphi,2)$ $=$ no-Salem closure & \F{} at $n=3$;
\C/\Pl{} general & Thm.~\ref{thm:parityrel}, \S\ref{sec:locking}\\
\hline
\end{tabular}
\end{center}

\section{The quadratic case of \textup{(EG)}, self-contained}
\label{app:quadEG}

For completeness we reproduce the elementary argument behind the \F{} portion
of \textup{(EG)}: \emph{no charge-admissible object of degree at most $2$ has
Mahler measure in $(1,\varphi)$.}

Degree-$1$ objects and reducible quadratics have nonzero integer roots, so
$\Mah\in\{1\}\cup\Z_{\ge2}$. Let $p=x^2-tx+n$ be irreducible, $t,n\in\Z$,
$n\neq0$.

\emph{Complex pair.} Then $n=\alpha\overline{\alpha}=|\alpha|^2\ge1$.
Admissibility puts $\alpha$ on a rational ray; if $n=1$ then $\alpha$ is
unimodular with rational angle, hence a root of unity and $\Mah=1$; if
$n\ge2$ then $\Mah=n\ge2$.

\emph{Real pair, $|n|\ge2$.} Real conjugates are automatically admissible
(angles in $\{0,\tfrac12\}$), and not both roots can lie in the closed unit
disk (their product has modulus $\ge2$). If both lie outside, $\Mah=|n|\ge2$;
if one lies strictly inside, $\Mah=|n|/|\text{inner root}|>|n|\ge2$.

\emph{Real units, $n=-1$.} The roots are $\frac{t\pm\sqrt{t^2+4}}{2}$; $t=0$
gives the reducible $x^2-1$, and for $|t|\ge1$,
\[
\Mah=\frac{|t|+\sqrt{t^2+4}}{2}\ \ge\ \frac{1+\sqrt5}{2}=\varphi,
\]
with equality exactly at $|t|=1$ (ledger entry~R): both gaps $|t|-1\ge0$ and
$\sqrt{t^2+4}-\sqrt5\ge0$ are nonnegative.

\emph{Real units, $n=+1$.} Real roots force $t^2>4$, i.e.\ $|t|\ge3$, and
$\Mah=\frac{|t|+\sqrt{t^2-4}}{2}\ge\frac{3+\sqrt5}{2}=\varphi^2>\varphi$
(ledger entry~R); $|t|\le2$ gives only cyclotomic or reducible cases.

Hence every quadratic value of $\Mah$ lies in $\{1\}\cup[\varphi,\infty)$,
with $\varphi$ attained by $x^2-x-1$. Beyond degree $2$, \textup{(EG)}
remains \C/\Pl{} as recorded in Remark~\ref{rem:gapstatus}; the structural
reason the gap is expected to survive --- the would-be occupants of
$(1,\varphi)$ near rational-angle lattices are Salem/Pisot numbers with
irrational angles, which are inadmissible --- is \cite[Lem.~6.1,
Thm.~6.2]{CMC}.

\begin{thebibliography}{9}

\bibitem{CMC} \emph{The Charge--Measure Coupling on a Spectral Semiring:
Cyclicity, Composition, and a Parity-Graded Mahler Floor},
math-research-pipelines technical note, 30 June 2026. Archived within the project repository deposit at DOI
\url{10.5281/zenodo.21121863} (v1.0.0, as
\texttt{papers/charge\_measure\_coupling.tex}); served at
\texttt{https://echo-s-studios.github.io/math-research-pipelines/}.
Appendix~\ref{app:inherited} lists every statement imported from it.

\bibitem{Kronecker} L.~Kronecker, \emph{Zwei S\"atze \"uber Gleichungen mit
ganzzahligen Coefficienten}, J.~Reine Angew.~Math.\ 53 (1857), 173--175.

\bibitem{Lehmer} D.~H.~Lehmer, \emph{Factorization of certain cyclotomic
functions}, Ann.\ of Math.\ 34 (1933), 461--479.

\bibitem{Salem} R.~Salem, \emph{Power series with integral coefficients},
Duke Math.~J.\ 12 (1945), 153--172.

\bibitem{Mahler} K.~Mahler, \emph{On some inequalities for polynomials in
several variables}, J.~London Math.~Soc.\ 37 (1962), 341--344.

\bibitem{Smyth} C.~J.~Smyth, \emph{On the product of the conjugates outside the
unit circle of an algebraic integer}, Bull.~London Math.~Soc.\ 3 (1971),
169--175.

\bibitem{Boyd} D.~W.~Boyd, \emph{Speculations concerning the range of Mahler's
measure}, Canad.~Math.~Bull.\ 24 (1981), 453--469.

\bibitem{Bertin} M.~J.~Bertin et al., \emph{Pisot and Salem Numbers},
Birkh\"auser, 1992.

\bibitem{Boyd1977} D.~W.~Boyd, \emph{Small Salem numbers}, Duke Math.~J.\
44 (1977), 315--328.

\bibitem{Dobrowolski} E.~Dobrowolski, \emph{On a question of Lehmer and the
number of irreducible factors of a polynomial}, Acta Arith.\ 34 (1979),
391--401.

\bibitem{Smyth1986} C.~J.~Smyth, \emph{Additive and multiplicative
relations connecting conjugate algebraic numbers}, J.~Number Theory 23
(1986), 243--254.

\bibitem{Mossinghoff} M.~J.~Mossinghoff, \emph{Polynomials with small
Mahler measure}, Math.~Comp.\ 67 (1998), 1697--1705.

\bibitem{EPSW} G.~Everest, A.~van der Poorten, I.~Shparlinski, and
T.~Ward, \emph{Recurrence Sequences}, Math.\ Surveys Monogr.\ 104,
Amer.\ Math.\ Soc., 2003.

\end{thebibliography}

\medskip
\noindent\textbf{Citation provenance:} reference~\cite{CMC} is the
companion note of this project (DOI \url{10.5281/zenodo.21121863}). The bibliographic metadata of
\cite{Kronecker}--\cite{Bertin} was machine-checked against the Crossref
registry on 1~July~2026. Confirmed in full: \cite{Lehmer}
(DOI~\texttt{10.2307/1968172}), \cite{Mahler}
(\texttt{10.1112/jlms/s1-37.1.341}), \cite{Smyth}
(\texttt{10.1112/blms/3.2.169}), and \cite{Bertin}
(\texttt{10.1007/978-3-0348-8632-1}, Birkh\"auser, 1992). Confirmed with
registry quirks noted: \cite{Kronecker} --- title, year 1857, and
pp.~173--175 confirmed; Crossref registers old Crelle volumes by year, while
the DOI \texttt{10.1515/crll.1857.53.173} encodes the traditional volume~53
used here. \cite{Boyd} --- title, volume~24, and pp.~453--469 confirmed at
DOI~\texttt{10.4153/cmb-1981-069-5}, whose registry issued-year field reads
1980 while the DOI and the standard citation year are 1981, retained here.
\cite{Salem} --- title, year 1945, and volume~12 confirmed at
DOI~\texttt{10.1215/s0012-7094-45-01213-0}; the registry record carries no
page field, so the range 153--172 stands as supplied. References
\cite{Boyd1977}, \cite{Dobrowolski}, \cite{Smyth1986}, and
\cite{Mossinghoff}, added at review, were machine-checked likewise, at DOIs
\url{10.1215/s0012-7094-77-04413-1} (Duke's registry record, as for
\cite{Salem}, carries no page field, so 315--328 stands as supplied),
\url{10.4064/aa-34-4-391-401}, \url{10.1016/0022-314x(86)90094-6}, and
\url{10.1090/s0025-5718-98-01006-0}. Reference~\cite{EPSW} was
machine-checked at DOI \url{10.1090/surv/104}.

\end{document}
